\documentclass{article}
\usepackage{amsmath}
\usepackage{amsfonts}
\usepackage{amssymb}
\usepackage{amsthm}
\usepackage{geometry}
\usepackage{algorithm}
\usepackage{algpseudocode}
\usepackage{minted}
\usepackage{graphicx}
\usepackage{float}

\geometry{a4paper, left=2cm, right=2cm, top=2cm, bottom=2cm}

\title{Image Processing: Homework 1}
\author{FU, YONG-WEI (112550107)}
\date{\today}

\begin{document}

\maketitle

\section{Question 1: Image Capture and Histograms}
\subsection{Composition Reasoning}
For this assignment, I chose to take the picture of the exterior structure of the 
staircase next to the library. I choose this place because I think that by changing 
the perspectives, I can show two different light intensity on the same scene. One side 
of it reflects the sunlight pretty well, which easily leads the overexposure as the 
requirement asks, and the other side of it is in the shadow, which can be underexposed, 
also meeting the requirement. But still, to meet the requirements, two photos were 
intentionally overexposed to blow out the sky and bright surfaces into pure white, 
while two were underexposed to crush the shadows beneath the structure into pure black.

\subsection{Histograms and Results}
\begin{figure}[H]
    \centering
    \includegraphics[width=\textwidth]{Q1_Base_Images/Q1_Histograms_Result.png}
    \caption{Four captured images converted to grayscale with their respective histograms.}
\end{figure}

The code for obtaining the histograms can be found in the \texttt{src} folder. Please refer 
to \texttt{q1\_histogram.py} for details.

\section{Question 2: Noise Generation}
\subsection{Histogram Analysis}
I choosed the second darker image to work on, as I found that it has a 
region that is in the shadows of the structure, which is all black.

\begin{figure}[H]
    \centering
    \includegraphics[width=\textwidth]{Q2_Noise_Results/Q2_Results_Plot.png}
    \caption{Full images with local bounding boxes (top) and local histograms of the selected smooth region (bottom).}
\end{figure}

After applying Gaussian noise to the image, in the selected dark region, a 
significant stacking effect is observed at the lower bound (pixel intensity 0). 
In the original local histogram, the peak is approximately 5000 pixels, at pixel value 0. 
After adding Gaussian noise, the peak height increases drastically (to roughly 14,000 for Var=100, and 17,500 for Var=900). 
This occurs because the initial pixel values are already ``near zero''; when the gaussian distribution applies 
negative shifts, the values underflow and are clipped at ``exactly zero''. Additionally, 
as variance increases, the right-side tail of the gaussian distribution becomes noticeably 
wider, therefore leads to more pixels being shifted to the higher pixel values.

And after applying Salt-and-Pepper noise, the histogram shows a two-peak distribution at 
the extreme pixel values (0 and 255). The peak at pixel value 0 increases significantly also, 
as more pixels are randomly set to black. Simultaneously, a new peak emerges at pixel value 255, 
as some pixels are randomly set to white. The height of both peaks increases with higher noise probability, 
indicating a bigger impact of the noise on the image.

The code for generating the noise and plotting the histograms can be found in the \texttt{src} 
folder. Please refer to \texttt{q2\_noise.py} for details.

\section{Question 3: Power-law (Gamma) Transformation}
\subsection{Contrast Correction}
To correct the awful contrast we have now, a Power-law transformation 
($s = c \cdot r^\gamma$) was applied. For the dark image, a fractional exponent 
($\gamma = 0.4$) was used to lift the dark regions up and reveal shadow details. 
For the bright image, a larger exponent ($\gamma = 2.5$) was applied to compress 
the highlights and restore depth.

\begin{figure}[H]
    \centering
    \includegraphics[width=\textwidth]{Q3_Contrast_Results/Q3_Powerlaw_Plot.png}
    \caption{Contrast correction using Power-law transformation.}
\end{figure}

The code for applying the Power-law transformation can be found in the 
\texttt{src} folder. Please refer to \texttt{q3\_powerlaw.py} for details.

\section{Question 4: Creative Edit}
\subsection{Choice of tool}
For this creative edit, I still chose to use Python along with the OpenCV and NumPy 
libraries instead of a traditional photo editing software. Python has a lot of 
open-source libraries that make image processing tasks straightforward and efficient, 
and including that I am more comfortable with programming approaches to image manipulation, and 
also the explanation for the implementation steps can be more precise and detailed, python is a 
perfect choice for this task. 

\subsection{Implementation and Steps}
This composite combines the contrast enhancement technique from Question 3 with 
image matting to create a surreal effect where the building appears to be burning from the deep shadows. 
As everyone knows that, the 24k library is the place where students study hard, and I always 
think that for a place like that, the energy around that place should be pretty intense, therefore the idea 
of fire just came into my mind, thus I decided to create this composite by blending them together.
The exact steps are as follows:

\begin{enumerate}
    \item \textbf{Power-law Transformation:} First, I took the original underexposed 
    image and applied a Power-law transformation with a fractional exponent ($\gamma = 0.4$). 
    This lifts the dark regions and recovers the original structure so the dark details of the image 
    is actually visible.
    \item \textbf{Shadow Masking:} Next, the contrast-corrected image was converted to grayscale. I 
    applied binary thresholding (with a threshold value of 80) to isolate the remaining extreme dark regions. 
    This step created a precise binary mask of the shadows.
    \item \textbf{Image Composition:} Finally, I used bitwise operations 
    (\texttt{cv2.bitwise\_not} and \texttt{cv2.bitwise\_and}) to extract a 
    downloaded fire background using the shadow mask. Using \texttt{cv2.add}, 
    I seamlessly composited the flames exclusively into the dark regions of the 
    brightened image.
\end{enumerate}

\begin{figure}[H]
    \centering
    \includegraphics[width=\textwidth]{Q4_Creative_Edit/Q4_Creative_Plot.png}
    \caption{Original dark image, Power-law correction, and final fire composite.}
\end{figure}

The complete code for this creative edit and image composition can be found in the \texttt{src} folder. Please refer to \texttt{q4\_creative.py} for details.

\end{document}